\chapter{Paths and Fundamental Groups}
\label{ch:viii09-paths-fundamental-groups}

The first genuinely algebraic invariant of a space is built from loops.  Paths
record how points are connected, while loops based at a chosen point can be
concatenated.  Passing to homotopy classes converts this concatenation into a
group law.  The resulting fundamental group detects holes, organizes covering
spaces, and encodes two-cell attaching loops in low-dimensional CW complexes.

\section{Paths}

\begin{definition}[Path]
\label{def:viii09-path}
A path in \(X\) from \(x_0\) to \(x_1\) is a continuous map
\[
\alpha:I\to X
\]
with
\[
\alpha(0)=x_0,
\qquad
\alpha(1)=x_1.
\]
\end{definition}

A loop based at \(x_0\) is a path satisfying
\[
\alpha(0)=\alpha(1)=x_0.
\]

\section{Path reversal}

For a path \(\alpha\), define
\[
\overline\alpha(t)=\alpha(1-t).
\]
Then \(\overline\alpha\) runs from the endpoint of \(\alpha\) back to its
starting point.

\section{Concatenation}

If
\[
\alpha(1)=\beta(0),
\]
define
\[
(\alpha*\beta)(t)=
\begin{cases}
\alpha(2t),&0\le t\le\frac12,\\
\beta(2t-1),&\frac12\le t\le1.
\end{cases}
\]

\begin{warning}[Convention]
\label{warn:viii09-convention}
This chapter reads \(\alpha*\beta\) as first traverse \(\alpha\), then
\(\beta\).  Some texts adopt the opposite multiplication convention.
\end{warning}

\section{Path homotopy}

\begin{definition}[Homotopy rel endpoints]
\label{def:viii09-path-homotopy}
Two paths \(\alpha,\beta:I\to X\) with the same endpoints are path-homotopic if
there exists
\[
H:I\times I\to X
\]
such that
\[
H(s,0)=\alpha(s),
\qquad
H(s,1)=\beta(s),
\]
and
\[
H(0,t)=x_0,
\qquad
H(1,t)=x_1.
\]
\end{definition}

This is exactly homotopy relative to \(\{0,1\}\) from VIII/01.

\section{The fundamental group}

\begin{definition}[Fundamental group]
\label{def:viii09-pi1}
The fundamental group of \(X\) at \(x_0\) is
\[
\pi_1(X,x_0)
=
\{\text{loops at }x_0\}/\text{path homotopy},
\]
with multiplication
\[
[\alpha][\beta]=[\alpha*\beta].
\]
\end{definition}

\begin{theorem}[Group structure]
\label{thm:viii09-group}
Path concatenation induces a well-defined group structure on
\(\pi_1(X,x_0)\).
\end{theorem}

\section{Identity and inverse}

The identity element is represented by the constant loop
\[
c_{x_0}(t)=x_0.
\]
The inverse of \([\alpha]\) is
\[
[\overline\alpha].
\]

Associativity of literal path concatenation holds only up to reparameterization,
but that is enough because reparameterized paths are homotopic rel endpoints.

\section{Well-defined multiplication}

\begin{proposition}[Compatibility with path homotopy]
\label{prop:viii09-well-defined}
If
\[
\alpha\simeq\alpha'
\quad\text{and}\quad
\beta\simeq\beta'
\]
rel endpoints, then
\[
\alpha*\beta\simeq\alpha'*\beta'
\]
rel endpoints.
\end{proposition}

Hence multiplication depends only on homotopy classes.

\section{Simply connected spaces}

\begin{definition}[Simply connected]
\label{def:viii09-simply-connected}
A space \(X\) is simply connected if it is path connected and
\[
\pi_1(X,x_0)=0
\]
for one, equivalently every, basepoint.
\end{definition}

Contractible spaces are simply connected.

\begin{example}
For \(n\ge2\),
\[
\pi_1(S^n)=0.
\]
\end{example}

\section{The circle}

\begin{theorem}[Fundamental group of the circle]
\label{thm:viii09-circle}
\[
\pi_1(S^1,1)\cong\mathbb Z.
\]
\end{theorem}

The integer records winding number.  The loop
\[
\omega_n(t)=e^{2\pi int}
\]
represents the integer \(n\).

\section{Winding and degree}

A loop in \(S^1\) is equivalently a based map
\[
S^1\to S^1.
\]
Its fundamental-group class is measured by the same integer that appeared as
degree in VIII/03.

\begin{proposition}[Degree and \(\pi_1(S^1)\)]
\label{prop:viii09-degree}
Under the identification \(\pi_1(S^1)\cong\mathbb Z\),
\[
[\omega_n]\longleftrightarrow n.
\]
\end{proposition}

\section{Basepoint change}

Let \(\gamma\) be a path from \(x_0\) to \(x_1\).

\begin{theorem}[Change of basepoint]
\label{thm:viii09-basepoint}
There is an isomorphism
\[
\gamma_\#:\pi_1(X,x_0)\to\pi_1(X,x_1)
\]
given by
\[
[\alpha]
\longmapsto
[\overline\gamma*\alpha*\gamma]
\]
with the concatenation convention used in this chapter, after interpreting the
order consistently.
\end{theorem}

Different choices of connecting path may change the isomorphism by an inner
automorphism.

\section{Induced homomorphisms}

For a based map
\[
f:(X,x_0)\to(Y,y_0),
\]
define
\[
f_*:\pi_1(X,x_0)\to\pi_1(Y,y_0)
\]
by
\[
f_*([\alpha])=[f\circ\alpha].
\]

\begin{theorem}[Functoriality]
\label{thm:viii09-functorial}
\[
(\operatorname{id}_X)_*=\operatorname{id}_{\pi_1(X)},
\qquad
(g\circ f)_*=g_*\circ f_*.
\]
\end{theorem}

\section{Homotopy invariance}

\begin{theorem}[Based homotopy invariance]
\label{thm:viii09-based-homotopy}
Based-homotopic maps induce the same homomorphism on fundamental groups.
\end{theorem}

\begin{corollary}[Homotopy equivalence]
\label{cor:viii09-equivalence}
A homotopy equivalence of path-connected spaces induces an isomorphism of
fundamental groups, after the usual basepoint identifications.
\end{corollary}

\section{Products}

\begin{theorem}[Product formula]
\label{thm:viii09-product}
For path-connected based spaces,
\[
\pi_1(X\times Y,(x_0,y_0))
\cong
\pi_1(X,x_0)\times\pi_1(Y,y_0).
\]
\end{theorem}

Hence
\[
\pi_1(T^2)\cong\mathbb Z^2.
\]

\section{Wedges and free groups: preview}

For a wedge of \(r\) circles,
\[
\bigvee_{j=1}^rS^1,
\]
the fundamental group is the free group
\[
F_r.
\]
A complete computation uses van Kampen's theorem, developed later in the
fundamental-group arc.

\section{Attaching loops and relations}

A two-cell attaching map
\[
\varphi:S^1\to X^1
\]
represents a loop class in
\[
\pi_1(X^1).
\]
Attaching the two-cell makes that loop null-homotopic in the enlarged space.
Thus two-cell attachments impose relations on the fundamental group.

\begin{example}[Torus relation]
\label{ex:viii09-torus-relation}
The torus one-skeleton has two loop generators \(a,b\).  Its two-cell attaches
along
\[
aba^{-1}b^{-1},
\]
forcing the commutator to become trivial.  This is the geometric source of
\[
\pi_1(T^2)\cong\mathbb Z^2.
\]
\end{example}

\section{Fundamental group as a homotopy invariant}

If two path-connected spaces have nonisomorphic fundamental groups, then they
cannot be homotopy equivalent.  In particular,
\[
S^1\not\simeq\{*\}
\]
because
\[
\pi_1(S^1)\cong\mathbb Z
\]
whereas the point has trivial fundamental group.

\section{Bridge to covering spaces}

The circle computation is most naturally understood using the covering map
\[
p:\mathbb R\to S^1,
\qquad
p(t)=e^{2\pi it}.
\]
Loops lift to paths in \(\mathbb R\), and the endpoint difference is an integer.
VIII/10 develops covering spaces systematically and turns this observation into
a general lifting theory.


\section{Exercises}

\begin{exercise}\label{exr:viii09-01}
Define a path from \(x_0\) to \(x_1\).
\end{exercise}
\begin{hint}
Use a continuous map on \(I\).
\end{hint}
\begin{solution}
A continuous \(\alpha:I\to X\) with \(\alpha(0)=x_0\) and \(\alpha(1)=x_1\).
\end{solution}

\begin{exercise}\label{exr:viii09-02}
Define a loop at \(x_0\).
\end{exercise}
\begin{hint}
Make the endpoints equal.
\end{hint}
\begin{solution}
A path \(\alpha\) with \(\alpha(0)=\alpha(1)=x_0\).
\end{solution}

\begin{exercise}\label{exr:viii09-03}
Define the reverse path \(\overline\alpha\).
\end{exercise}
\begin{hint}
Reverse the parameter.
\end{hint}
\begin{solution}
\(\overline\alpha(t)=\alpha(1-t)\).
\end{solution}

\begin{exercise}\label{exr:viii09-04}
When is \(\alpha*\beta\) defined?
\end{exercise}
\begin{hint}
Match the endpoint of the first to the start of the second.
\end{hint}
\begin{solution}
When \(\alpha(1)=\beta(0)\).
\end{solution}

\begin{exercise}\label{exr:viii09-05}
Write the concatenation formula.
\end{exercise}
\begin{hint}
Give each path half the interval.
\end{hint}
\begin{solution}
\((\alpha*\beta)(t)=\alpha(2t)\) for \(t\le1/2\), and \(\beta(2t-1)\) for \(t\ge1/2\).
\end{solution}

\begin{exercise}\label{exr:viii09-06}
What must a path homotopy keep fixed?
\end{exercise}
\begin{hint}
Use homotopy rel endpoints.
\end{hint}
\begin{solution}
The initial and terminal endpoints.
\end{solution}

\begin{exercise}\label{exr:viii09-07}
Define \(\pi_1(X,x_0)\).
\end{exercise}
\begin{hint}
Take based loops modulo endpoint-fixing homotopy.
\end{hint}
\begin{solution}
The set of based loop homotopy classes with concatenation.
\end{solution}

\begin{exercise}\label{exr:viii09-08}
What is the identity in \(\pi_1(X,x_0)\)?
\end{exercise}
\begin{hint}
Use the constant loop.
\end{hint}
\begin{solution}
\([c_{x_0}]\).
\end{solution}

\begin{exercise}\label{exr:viii09-09}
What is the inverse of \([\alpha]\)?
\end{exercise}
\begin{hint}
Reverse the loop.
\end{hint}
\begin{solution}
\([\overline\alpha]\).
\end{solution}

\begin{exercise}\label{exr:viii09-10}
Why is multiplication well-defined?
\end{exercise}
\begin{hint}
Concatenate representatives of homotopies.
\end{hint}
\begin{solution}
Path-homotopic representatives have path-homotopic concatenations.
\end{solution}

\begin{exercise}\label{exr:viii09-11}
Is literal path concatenation strictly associative?
\end{exercise}
\begin{hint}
Compare parametrizations.
\end{hint}
\begin{solution}
No; it is associative up to endpoint-fixing reparameterization homotopy.
\end{solution}

\begin{exercise}\label{exr:viii09-12}
Define simply connected.
\end{exercise}
\begin{hint}
Require connectivity and trivial \(\pi_1\).
\end{hint}
\begin{solution}
Path connected with trivial fundamental group.
\end{solution}

\begin{exercise}\label{exr:viii09-13}
What is \(\pi_1(S^n)\) for \(n\ge2\)?
\end{exercise}
\begin{hint}
Higher spheres are simply connected.
\end{hint}
\begin{solution}
Trivial.
\end{solution}

\begin{exercise}\label{exr:viii09-14}
What is \(\pi_1(S^1)\)?
\end{exercise}
\begin{hint}
Use winding number.
\end{hint}
\begin{solution}
\(\mathbb Z\).
\end{solution}

\begin{exercise}\label{exr:viii09-15}
What integer does \(\omega_n(t)=e^{2\pi int}\) represent?
\end{exercise}
\begin{hint}
Count windings.
\end{hint}
\begin{solution}
\(n\).
\end{solution}

\begin{exercise}\label{exr:viii09-16}
How is degree related to \(\pi_1(S^1)\)?
\end{exercise}
\begin{hint}
Both classify circle maps.
\end{hint}
\begin{solution}
The degree equals the winding-number integer under \(\pi_1(S^1)\cong\mathbb Z\).
\end{solution}

\begin{exercise}\label{exr:viii09-17}
What does a based map \(f:X\to Y\) induce?
\end{exercise}
\begin{hint}
Compose loops with \(f\).
\end{hint}
\begin{solution}
A homomorphism \(f_*:\pi_1(X)\to\pi_1(Y)\).
\end{solution}

\begin{exercise}\label{exr:viii09-18}
State functoriality for \(g\circ f\).
\end{exercise}
\begin{hint}
Induced maps compose.
\end{hint}
\begin{solution}
\((g\circ f)_*=g_*\circ f_*\).
\end{solution}

\begin{exercise}\label{exr:viii09-19}
What do based-homotopic maps induce?
\end{exercise}
\begin{hint}
Use homotopy invariance.
\end{hint}
\begin{solution}
The same homomorphism on \(\pi_1\).
\end{solution}

\begin{exercise}\label{exr:viii09-20}
What does a homotopy equivalence induce on \(\pi_1\)?
\end{exercise}
\begin{hint}
Use a homotopy inverse.
\end{hint}
\begin{solution}
An isomorphism, after appropriate basepoint identification.
\end{solution}

\begin{exercise}\label{exr:viii09-21}
State the product formula.
\end{exercise}
\begin{hint}
Project loops to both factors.
\end{hint}
\begin{solution}
\(\pi_1(X\times Y)\cong\pi_1(X)\times\pi_1(Y)\).
\end{solution}

\begin{exercise}\label{exr:viii09-22}
Compute \(\pi_1(T^2)\).
\end{exercise}
\begin{hint}
Use \(T^2=S^1\times S^1\).
\end{hint}
\begin{solution}
\(\mathbb Z^2\).
\end{solution}

\begin{exercise}\label{exr:viii09-23}
What does attaching a two-cell along a loop do to that loop class?
\end{exercise}
\begin{hint}
The disk fills the loop.
\end{hint}
\begin{solution}
It makes the attaching loop null-homotopic in the enlarged space.
\end{solution}

\begin{exercise}\label{exr:viii09-24}
What is the bridge to VIII/10?
\end{exercise}
\begin{hint}
Use the universal cover of the circle.
\end{hint}
\begin{solution}
Covering-space lifting turns loop classes into endpoint data of lifted paths.
\end{solution}


\section{Solved problem dossiers}

\begin{problem}[Reverse twice]\label{prob:viii09-01}
Show \(\overline{\overline\alpha}=\alpha\).
\end{problem}
\begin{solution}
By definition, \(\overline{\overline\alpha}(t)=\overline\alpha(1-t)=\alpha(1-(1-t))=\alpha(t)\).
\end{solution}

\begin{problem}[Concatenation endpoints]\label{prob:viii09-02}
Verify the endpoints of \(\alpha*\beta\).
\end{problem}
\begin{solution}
At \(t=0\) the value is \(\alpha(0)\); at \(t=1\) it is \(\beta(1)\).  At \(t=1/2\), both formulas equal the common point \(\alpha(1)=\beta(0)\).
\end{solution}

\begin{problem}[Identity loop]\label{prob:viii09-03}
Why is \(\alpha*c_{x_1}\) homotopic rel endpoints to \(\alpha\)?
\end{problem}
\begin{solution}
Reparameterize the traversal of \(\alpha\) so that it reaches the endpoint later and removes the terminal pause.  Endpoint-preserving reparameterization gives the required homotopy.
\end{solution}

\begin{problem}[Inverse loop]\label{prob:viii09-04}
Why is \(\alpha*\overline\alpha\) null-homotopic rel basepoint?
\end{problem}
\begin{solution}
Shrink the out-and-back traversal symmetrically toward the basepoint; at each stage traverse only an initial segment of \(\alpha\) and immediately return along it.
\end{solution}

\begin{problem}[Associativity]\label{prob:viii09-05}
Explain why \((\alpha*\beta)*\gamma\) and \(\alpha*(\beta*\gamma)\) represent the same fundamental-group element.
\end{problem}
\begin{solution}
They traverse the same three paths in the same order with different time allocations.  A continuous reparameterization homotopy rel endpoints changes one allocation into the other.
\end{solution}

\begin{problem}[Circle generator]\label{prob:viii09-06}
Show that \(\omega_1(t)=e^{2\pi it}\) generates \(\pi_1(S^1)\cong\mathbb Z\).
\end{problem}
\begin{solution}
Every loop class has an integer winding number \(n\), and \(\omega_n\) is homotopic to the \(n\)-fold concatenation of \(\omega_1\) for \(n>0\), with reversed copies for \(n<0\).
\end{solution}

\begin{problem}[Degree and winding]\label{prob:viii09-07}
For \(f(z)=z^m\), compute the induced map \(f_*:\pi_1(S^1)\to\pi_1(S^1)\).
\end{problem}
\begin{solution}
The generator loop \(\omega_1\) is sent to \(t\mapsto e^{2\pi imt}=\omega_m\).  Thus \(f_*\) is multiplication by \(m\).
\end{solution}

\begin{problem}[Contractible target]\label{prob:viii09-08}
Show that a contractible path-connected space has trivial fundamental group.
\end{problem}
\begin{solution}
The identity is homotopic to a constant map.  Homotopy invariance makes the identity homomorphism on \(\pi_1\) equal to the zero homomorphism induced by the constant map, so the group is trivial.
\end{solution}

\begin{problem}[Product loop]\label{prob:viii09-09}
Describe a loop in \(X\times Y\) in terms of loops in the factors.
\end{problem}
\begin{solution}
A loop \(\alpha(t)=(\alpha_X(t),\alpha_Y(t))\) is exactly a pair of loops.  Homotopies and concatenation occur componentwise, yielding the direct-product formula.
\end{solution}

\begin{problem}[Torus group]\label{prob:viii09-10}
Compute \(\pi_1(T^2)\) from the product formula.
\end{problem}
\begin{solution}
\(\pi_1(S^1\times S^1)\cong\mathbb Z\times\mathbb Z\).
\end{solution}

\begin{problem}[Basepoint path]\label{prob:viii09-11}
Explain why path-connected spaces have isomorphic fundamental groups at different basepoints.
\end{problem}
\begin{solution}
Choose a path joining the basepoints and conjugate loops by that path and its reverse.  This gives an isomorphism, with different path choices differing by inner automorphisms.
\end{solution}

\begin{problem}[Homotopy-equivalence obstruction]\label{prob:viii09-12}
Why can \(S^1\) not be homotopy equivalent to a point?
\end{problem}
\begin{solution}
Homotopy equivalence would induce an isomorphism of fundamental groups, but \(\pi_1(S^1)\cong\mathbb Z\) and the point has trivial fundamental group.
\end{solution}

\begin{problem}[Higher sphere comparison]\label{prob:viii09-13}
Use \(\pi_1\) to distinguish \(S^1\) from \(S^2\).
\end{problem}
\begin{solution}
\(\pi_1(S^1)\cong\mathbb Z\) while \(\pi_1(S^2)=0\), so the spaces are not homotopy equivalent.
\end{solution}

\begin{problem}[Wedge preview]\label{prob:viii09-14}
What group should be expected for \(S^1\vee S^1\)?
\end{problem}
\begin{solution}
The expected group is the free group \(F_2\) on the two circle loops.  Van Kampen's theorem gives the formal computation.
\end{solution}

\begin{problem}[Torus attaching relation]\label{prob:viii09-15}
Explain geometrically why the torus two-cell kills the commutator.
\end{problem}
\begin{solution}
The boundary of the two-cell is the loop \(aba^{-1}b^{-1}\).  Once the disk is attached, that boundary loop bounds the disk and is therefore null-homotopic, imposing \([a,b]=1\).
\end{solution}

\begin{problem}[Null two-cell attachment]\label{prob:viii09-16}
If a two-cell attaches along a null-homotopic loop in \(X\), what does VIII/08 predict?
\end{problem}
\begin{solution}
The result is homotopy equivalent to \(X\vee S^2\), because the attaching loop can be replaced by the constant loop.
\end{solution}

\begin{problem}[Functoriality check]\label{prob:viii09-17}
Prove \((g\circ f)_*=g_*\circ f_*\) on a loop class.
\end{problem}
\begin{solution}
For \([\alpha]\), \((g\circ f)_*[\alpha]=[g\circ f\circ\alpha]=g_*([f\circ\alpha])=(g_*\circ f_*)[\alpha]\).
\end{solution}

\begin{problem}[Constant map]\label{prob:viii09-18}
What homomorphism on fundamental groups is induced by a constant based map?
\end{problem}
\begin{solution}
Every loop is sent to the constant loop, so the induced homomorphism is trivial.
\end{solution}

\begin{problem}[Circle covering preview]\label{prob:viii09-19}
Lift \(\omega_n(t)=e^{2\pi int}\) to a path in \(\mathbb R\) starting at \(0\).
\end{problem}
\begin{solution}
The lift is \(\widetilde\omega_n(t)=nt\), because \(e^{2\pi i(nt)}=\omega_n(t)\).  Its endpoint is the integer \(n\).
\end{solution}

\begin{problem}[Why coverings next?]\label{prob:viii09-20}
Explain how the circle cover suggests a general method for studying \(\pi_1\).
\end{problem}
\begin{solution}
Loop classes can be detected by how their lifts move between points in a fiber.  VIII/10 generalizes path lifting and endpoint behavior to arbitrary covering spaces.
\end{solution}

\begin{problem}[Loop reparameterization]\label{prob:viii09-21}
Show that changing the speed of a loop by an endpoint-preserving map \(\rho:I\to I\) does not change its fundamental-group class.
\end{problem}
\begin{solution}
The homotopy \(H(s,t)=\alpha((1-t)s+t\rho(s))\) works when the interpolation remains endpoint preserving; more generally VIII/01's reparameterization principle provides the endpoint-fixing homotopy.
\end{solution}

\begin{problem}[Path-connected basepoint independence]\label{prob:viii09-22}
Why is the statement '\(\pi_1(X)\)' usually acceptable without naming a basepoint for a path-connected space?
\end{problem}
\begin{solution}
All basepoint fundamental groups are isomorphic.  The isomorphism is not completely canonical in nonabelian groups, but the group isomorphism type is independent of basepoint.
\end{solution}

\begin{problem}[Two-cell relation principle]\label{prob:viii09-23}
State the conceptual relation between VIII/08 and VIII/09 for two-dimensional CW complexes.
\end{problem}
\begin{solution}
VIII/08 says only the homotopy class of the attaching loop matters.  VIII/09 packages based loop homotopy classes into \(\pi_1\), so attaching data become group elements and relations.
\end{solution}

\begin{problem}[Bridge dossier]\label{prob:viii09-24}
Describe the role of \(p:\mathbb R\to S^1\), \(p(t)=e^{2\pi it}\), in the next chapter.
\end{problem}
\begin{solution}
It is the model covering map.  Paths and loops lift uniquely after choosing a starting point, and a loop's winding number is recovered from the endpoint of its lift.
\end{solution}
